
%Les données publiées le 1\up{er } mars 2023 par le ministère de la transition écologique sur les immatriculations de véhicules particuliers en France en 2022 contiennent les informations suivantes :
%
%	\begin{itemize}
%		\item $22,86\,\%$ des véhicules étaient des véhicules neufs ;
%		\item $8,08\,\%$ des véhicules neufs étaient des hybrides rechargeables;
%		\item $1,27\,\%$ des véhicules d'occasion (c'est-à-dire qui ne sont pas neufs) étaient des hybrides rechargeables.
%	\end{itemize}
%
%	\emph{Dans tout l'exercice, les probabilités seront arrondies au dix-millième.}
\textbf{Partie A}

\medskip

%Dans cette partie, on considère un véhicule particulier immatriculé en France en 2022.
%
%On note :
%
%\begin{itemize}
%	\item $N$ l'évènement \og le véhicule est neuf \fg{};
%	\item $R$ l'évènement \og le véhicule est hybride rechargeable\fg{};
%	\item $\overline{N}$ et $\overline{R}$ les évènements contraires des évènements contraires de $N$ et $R$.
%\end{itemize}

\begin{enumerate}
\item \phantom{21}  %Représenter la situation par un arbre pondéré.Arbre pondéré de probilités :
\begin{center}
\pstree[treemode=R,levelsep = 3.5cm]{\TR{}}
{\pstree{\TR{$N~~$}\taput{\np{0,2286}}}
	{\TR{$R$}\taput{\np{0,0808}}
	\TR{$\overline{R}$}\tbput{\red 0,9192}
	}
\pstree{\TR{$\overline{N}~~$}\tbput{\red\np{0,7714}}}
	{\TR{$R$}\taput{\np{0,0127}}
	\TR{$\overline{R}$}\tbput{\red 0,9873}
	}
}
\end{center}

\item %Calculer la probabilité que ce véhicule soit neuf et hybride rechargeable.
On calcule $p(N \cap R) = p(N) \times p_N(R) = \np{0,2286} \times \np{0,0808} = \np{0,018471}$ soit \np{0,0185} à $10^{-4}$ près.
\item %Démontrer que la valeur arrondie au dix-millième de la probabilité que ce véhicule soit hybride rechargeable est \np{0,0283}.
On a de même $p\left(\overline{N} \cap R\right) = p\left(N\right) \times p_{\overline{N}}(R) = \np{0,7714} \times \np{0,0127} = \np{0,009797}$ soit \np{0,0098} à $10^{-4}$ près.

D'après la loi des probabilités totales :

$p(R) = p(N \cap R) + p\left(\overline{N} \cap R\right) \approx \np{0,0185} + \np{0,0098}$

$p(R) \approx \np{0,0283}$.

\item %Calculer la probabilité que ce véhicule soit neuf sachant qu'il est hybride rechargeable.
On a $p_R(N) = \dfrac{p(R \cap N)}{p(R)} = \dfrac{p(N \cap R)}{p(R)} \approx \dfrac{\np{0,0185}}{\np{0,0283}} \approx \np{0,65371}$ soit \np{0,6537} à $10^{-4}$ près.
\end{enumerate}

\textbf{Partie B}

\medskip

%Dans cette partie, on choisit $500$ véhicules particuliers hybrides rechargeables immatriculés en France en 2022.
%
%Dans la suite, on admettra que la probabilité qu'un tel véhicule soit neuf est égale à $0,65$.
%
%On assimile le choix de ces 500 véhicules à un tirage aléatoire avec remise.
%
%On appelle $X$ la variable aléatoire représentant le nombre de véhicules neufs parmi les $500$ véhicules choisis.

\medskip

\begin{enumerate}
\item %On admet que la variable aléatoire $X$ suit une loi binomiale. Préciser la valeur de ses paramètres.

$X$ suit donc la loi binomiale $\mathcal{B}(n = 500, \: p = 0,65)$.
\item %Déterminer la probabilité qu'exactement $325$ de ces véhicules soient neufs.
On calcule $\binom{500}{325} \times 0,65^{325} \times 0,35^{175} \approx \np{0,037384}$ soit \np{0,0374} à $10^{-4}$ près. (Utiliser la fonction binomiale de la calculatrice  si les capacités de calcul de celle-ci sont dépassées).
\item %Déterminer la probabilité $p(X \geqslant 325)$ puis interpréter le résultat dans le contexte de l'exercice.
On a $p(X \geqslant 325) = 1 - p(X \leqslant 324)$ soit d'après la calculatrice \np{0,47944}, donc 

$p(X \geqslant 325) \approx 1 - \np{0,4794} = \np{0,5206}$ à $10^{-4}$ près.
\end{enumerate}

\textbf{Partie C}

\medskip

%On choisit désormais $n$ véhicules particuliers hybrides rechargeables immatriculés en France en 2022, où $n$ désigne un entier naturel strictement positif.
%
%On rappelle que la probabilité qu'un tel véhicule soit neuf est égale à $0,65$.
%
%On assimile le choix de ces $n$ véhicules à un tirage aléatoire avec remise.

\medskip

\begin{enumerate}
\item %Donner l'expression en fonction de $n$ de la probabilité $p_{n}$ que tous ces véhicules soient d'occasion.
On a $p_n = (1 - 0,65)^n = 0,35^n$.

\item On a $q_n = 1 - p_n$.

On cherche donc $n$ tel que $1 - p_n \geqslant \np{0,9999} \iff p_n \leqslant \np{0,0001}$, soit par croissance de la fonction logarithme népérien :

$n\ln 0,35 \leqslant \ln 0\np{0,0001} \iff n \geqslant \dfrac{\ln \np{0,0001}}{\ln 0,35}$ car $\dfrac{1}{\ln 0,35} < 0$.

D'après la calculatrice $\dfrac{\ln \np{0,0001}}{\ln 0,35} \approx 8,8$.

Il faut prendre au minimum $n = 9$.


%On note $q_{n}$ la probabilité qu'au moins un de ces véhicules soit neuf. En résolvant une inéquation, déterminer la plus petite valeur de $n$ telle que $q_{n} \geqslant \np{0,9999}$.
\end{enumerate}

